% ECONOMICS_PROBLEM: {"id": "D47-540", "title": "Interacting platform designs", "jel_code": "D47", "source_id": "OP-0401"}
% !TeX program = xelatex
\documentclass[11pt,a4paper]{article}
\usepackage[margin=23mm]{geometry}
\usepackage{fontspec}
\setmainfont{Latin Modern Roman}
\usepackage{amsmath,amssymb,mathtools}
\usepackage{xcolor,enumitem,booktabs,longtable,array,xurl,hyperref}
\definecolor{muted}{HTML}{53636C}
\hypersetup{colorlinks=true,urlcolor=blue,linkcolor=blue}
\setlength{\parindent}{0pt}
\setlength{\parskip}{5pt}
\newcommand{\R}{\mathbb R}
\newcommand{\E}{\mathbb E}
\newcommand{\N}{\mathbb N}
\newcommand{\ind}{\mathbf 1}
\newcommand{\Var}{\operatorname{Var}}
\newcommand{\Cov}{\operatorname{Cov}}
\newcommand{\supp}{\operatorname{supp}}
\DeclareMathOperator*{\argmin}{arg\,min}
\DeclareMathOperator*{\argmax}{arg\,max}
\newcommand{\entrymeta}[3]{{\sffamily\small\color{muted}\textbf{\texttt{#1}}\quad|\quad #2\par #3\par}}
\newcommand{\Problemref}[1]{\texttt{#1}}
\newcommand{\JELref}[2]{\texttt{#2} (JEL \texttt{#1})}
\begin{document}
\section*{Interacting platform designs}
\textbf{Problem ID:} \texttt{D47-540}\quad \textbf{JEL:} \texttt{D47}\par
% BEGIN ECONOMICS STATEMENT
\label{problem:OP-0401}
{\sffamily\small\color{muted}Original subject group: Digital, platform and data economics\par}
\label{definitions:problem:30007683}
\textbf{Audit classification:} Published research agenda. The source explicitly motivates combinatorial interaction experiments. Eight fully observed cells are already identifiable; meaningful research must concern unobserved combinations, growing dimension or behavioral adjustment.
\subsection*{Definitions and precise research target}
Consider an online retail platform partitioned into geographically separate markets where cross-market spillovers are negligible. Let \(r,b,d\in\{0,1\}\) indicate respectively a price-sensitive ranking, a quality badge and standardized product disclosure. The eight triples \(z=(r,b,d)\) define feasible joint interventions; \(z=0\) means the all-zero baseline. Under each \(z\), sellers choose prices, quality and whether to participate after observing platform rules. Let \(W(z)\) be equilibrium total surplus per potential buyer, defined as consumer utility net of search and payments plus seller and platform profits net of real operating and quality costs. Transfers among these parties cancel. Define \(\Delta(z)=W(z)-W(0)\) and, for ranking and badge interaction with disclosure fixed at zero, \(I_{rb}=W(1,1,0)-W(1,0,0)-W(0,1,0)+W(0,0,0)\). Analogous contrasts describe the remaining pair and triple interactions. Identify these quantities using market-level factorial assignment with sufficient time for seller adjustment, or a validated structural counterfactual that permits those adjustments. Record both immediate behavior and subsequent price, quality and entry responses. Determine which sparse experimental designs recover the signs of economically material interactions under declared smoothness or sparsity restrictions. This is a joint-intervention welfare-design problem; additive individual-feature effects are an assumption to be tested, not the desired equilibrium estimand.

\textbf{Definitions, assumptions and mathematical qualifications.} An equilibrium-selection rule, identical shock distribution and common eligible-buyer denominator define \(W(z)\) for each of eight policies. The triple interaction is \(I_{rbd}=W(1,1,1)-W(1,1,0)-W(1,0,1)-W(0,1,1)+W(1,0,0)+W(0,1,0)+W(0,0,1)-W(0,0,0)\). Market randomization requires strictly positive assignment probability to each tested cell. With all eight cells observed, these finite contrasts are ordinary estimands; they are not an unsolved combinatorial theorem. In the \(K\)-feature extension, state \(W(z)=\sum_{S\subseteq\{1,\ldots,K\}}\beta_S\prod_{k\in S}z_k\), the maximum interaction order and an explicit sparsity restriction before choosing a reduced design.
\subsection*{Variants and assumption changes}
\begin{enumerate}
\item \textbf{Editorial, complete factorial.} Randomize all eight cells and estimate each pair and triple interaction with market-level uncertainty; no sparsity restriction is needed.
\item \textbf{Editorial, large feature set.} Let \(K\) grow, allow at most \(s\) nonzero interactions of order at most \(d\), and seek budget-limited designs that recover material interaction signs under stated signal and noise bounds.
\end{enumerate}
\subsection*{References and source audit}
\textbf{Checked source.} 2025 conference draft, Section 1.1, second research direction, PDF pp. 4--5. Full primary text retrieved. \url{https://conference.nber.org/confer/2025/DTs25/farronato.pdf}.
\begin{enumerate}\raggedright
\item\hypertarget{definitions:source:30007683:1}{} Chiara Farronato. 2025. \emph{Digital Platforms as Information Aggregators: Implications for their Design and Regulation}. 2025 conference draft, Section 1.1, second research direction, PDF pp. 4--5.
\url{https://conference.nber.org/confer/2025/DTs25/farronato.pdf}.
\end{enumerate}

\subsection*{Common conventions: Source mathematical conventions}

These conventions apply to each entry unless explicitly replaced.

\textbf{Models and identification.} A primitive model \(m\) specifies all probability laws, feasible choices, information, equilibrium and measurement rules used by its target. The admissible class \(\mathcal M\) is fixed independently of outcomes. If \(P_O(m)\) is its observable probability law and \(\tau(m)\) the target, its identified set at observable law \(P\) is
\[
 \mathcal I_\tau(P)=\{\tau(m):m\in\mathcal M,\ P_O(m)=P\}.
\]
Point identification means that this set is a singleton. An empty set signals incompatibility of data and assumptions. Coordinate bounds need not describe the joint set. Bounds are sharp only if every retained target is attainable or arbitrarily approachable by compatible models. Finite bounds require explicit restrictions on unobserved quantities; unrestricted confounding, hidden wealth, extrapolation or arbitrary equilibrium selection can yield uninformative sets.

\textbf{Causal interventions.} A potential outcome \(Y_i(p)\) is the outcome under a complete policy regime \(p\), including timing, financing, information and market spillovers. The target population and units are fixed before treatment. With interference, use the complete assignment vector or a justified exposure mapping; individual no-interference notation alone cannot identify an economy-wide intervention. A difference of mean potential outcomes does not require the cross-world joint distribution. Conditioning on post-treatment migration, survival, enrollment or employment changes the estimand and needs separate justification.

\textbf{Statistical statements.} An estimator, confidence set, forecast or test specifies a sampling law, model class and loss/coverage criterion. Uniform \((1-\alpha)\)-coverage means \(\inf_{m\in\mathcal M}\Pr_m\{\tau(m)\in C_n\}\geq1-\alpha\), or an explicitly asymptotic version. The whole parameter space trivially has coverage, so informativeness requires a stated width, risk or power objective. A forecasting improvement is relative to a named benchmark, information set and evaluation population.

\textbf{Equilibrium and optimization.} An equilibrium comprises feasible plans, optimality, beliefs where needed, budget constraints and all market/resource clearing. The primitives or a stated benchmark define each correspondence \(\mathcal E(p;m)\). Nonemptiness is assumed or proved, never inferred just from compact policy sets. A selected-equilibrium target uses a declared selection rule; a robust target quantifies over all permitted equilibria. Use suprema or infima unless attainment is established. An \(\varepsilon\)-optimal policy is feasible and within \(\varepsilon\) of the finite optimum value. Report an empty feasible set separately.

\textbf{Welfare and accounting.} Normative utility, interpersonal normalization, social weights, numeraire, discounting and the use of public receipts are primitives. Behavioral utility need not coincide with the welfare criterion. Household welfare includes ownership income and rents once; adding profits again to compensating variation generally double-counts them. A monetary equivalent is defined by an explicit transfer and price convention, with monotonicity and range conditions. Average effects, mechanism contributions, transfers and welfare losses are different targets. Infinite sums require convergence and dynamic feasible plans require terminal or no-Ponzi conditions.

\textbf{Source versions.} A working-paper date, revision date and journal publication year are distinguished. A DOI for a journal version is not automatically the DOI of an earlier manuscript. Locators refer to the stated version and printed pages where specified. A metadata-only check does not verify a claim about a full-text conclusion. Every entry's source subsection narrows the provenance claim accordingly.
% END ECONOMICS STATEMENT
\end{document}
